% !TEX root = ../main.tex
% ================== SOLUCAO PROPOSTA ==================
\secao{Solução proposta}

A câmara é uma caixa de acrílico com três lâmpadas W5W (12\,V/5\,W) no teto, uma chapa de
alumínio como massa térmica e o sensor LM35 suspenso no centro. O ESP32 lê a temperatura,
calcula a ação de controle PI e gera o PWM que chaveia as lâmpadas pelo MOSFET.

% ---------- Diagrama da malha de controle (desenhado em TikZ, nao precisa de imagem) ----------
\titulofig{Malha de controle da câmara térmica}
\begin{center}
\begin{tikzpicture}[
    font=\normalsize,
    bloco/.style={draw, thick, rounded corners=6pt, fill=azultitulo!8, align=center,
                  minimum height=3.8cm, minimum width=8.2cm},
    seta/.style={-{Latex[length=7mm]}, line width=2.5pt},
    node distance=2.2cm]
  \node (sp)   {Setpoint};
  \node[draw, thick, circle, minimum size=1.6cm, right=1.6cm of sp] (soma) {};
  \node at (soma) {\(\Sigma\)};
  \node[bloco, right=of soma] (pi)  {Controlador PI\\(anti-windup)};
  \node[bloco, right=of pi]   (pwm) {PWM LEDC\\500\,Hz / 10 bits};
  \node[bloco, below=2.6cm of pwm] (mos) {MOSFET IRLZ44N\\+ 3 lâmpadas W5W};
  \node[bloco, left=of mos]   (cam) {Câmara\\(acrílico + alumínio)};
  \node[bloco, left=of cam]   (sen) {LM35 + ADC1\\(média de 16)};
  \draw[seta] (sp) -- (soma);
  \draw[seta] (soma) -- node[above]{\(e\)} (pi);
  \draw[seta] (pi) -- node[above]{\(u\)} (pwm);
  \draw[seta] (pwm) -- (mos);
  \draw[seta] (mos) -- (cam);
  \draw[seta] (cam) -- node[above]{\(T\)} (sen);
  \draw[seta] (sen.west) -| node[pos=0.8, left]{\(-\)} (soma.south);
\end{tikzpicture}
\end{center}

\vspace{0.5cm}
O controlador PI calcula o duty cycle a partir do erro \(e = T_{sp} - T\):
\begin{equation*}
  u(t) = K_p\, e(t) + K_i \int e(t)\, dt
\end{equation*}
Como a potência nas lâmpadas cai com o quadrado da tensão, o duty é corrigido pela tensão da
bateria:
\begin{equation*}
  D_{corr} = D_{PI} \cdot \left(\frac{12}{V_{bat}}\right)^{2}
\end{equation*}

% ---------- Cadeia de potencia (TikZ) ----------
\titulofig{Cadeia de potência e proteções}
\begin{center}
\begin{tikzpicture}[
    font=\normalsize,
    b/.style={draw, thick, rounded corners=6pt, fill=cinzabox, align=center,
              minimum height=3.4cm, minimum width=6.6cm},
    seta/.style={-{Latex[length=6mm]}, line width=2pt},
    node distance=1.3cm]
  \node[b] (bat) {Pack 3S 18650\\+ BMS};
  \node[b, right=of bat] (d)   {Schottky\\SS34};
  \node[b, right=of d]   (f)   {Fusível\\3\,A + TVS};
  \node[b, right=of f]   (k)   {Termostato\\KSD9700};
  \node[b, below=1.6cm of k] (l) {Lâmpadas\\W5W};
  \node[b, below=1.6cm of d] (buck) {Buck\\12\,\(\to\)\,5\,V};
  \node[b, left=of buck] (esp) {ESP32\\DevKit v1};
  \draw[seta] (bat) -- (d);
  \draw[seta] (d) -- (f);
  \draw[seta] (f) -- (k);
  \draw[seta] (k) -- (l);
  \draw[seta] (d) -- (buck);
  \draw[seta] (buck) -- (esp);
\end{tikzpicture}
\end{center}

\titulofig{Protótipo da câmara térmica}
\imagem[0.85\linewidth][15cm]{figuras/prototipo.jpg}
